% Generated draft: requires mathematical and source review.
\documentclass[11pt]{article}
\usepackage[margin=25mm]{geometry}
\usepackage{fontspec}
\setmainfont{Latin Modern Roman}
\usepackage{amsmath,amssymb,mathtools,unicode-math}
\setmathfont{Latin Modern Math}
\usepackage{enumitem,needspace,xurl,hyperref}
\setlength{\parindent}{0pt}
\setlength{\parskip}{5pt}
\setlength{\emergencystretch}{3em}
\newcommand{\N}{\mathbb{N}}
\newcommand{\Z}{\mathbb{Z}}
\newcommand{\Q}{\mathbb{Q}}
\newcommand{\R}{\mathbb{R}}
\newcommand{\C}{\mathbb{C}}
\newcommand{\F}{\mathbb{F}}
\newcommand{\A}{\mathbb{A}}
\newcommand{\PP}{\mathbb{P}}
\newcommand{\E}{\mathbb{E}}
\newcommand{\eps}{\varepsilon}
\newcommand{\dd}{\,\mathrm{d}}
\begin{document}
\section*{UnsolvedMath ID: 801126}
\textbf{Source problem code:} catalogue entry
\subsection{Definitions and mathematical statement}
The supplied formulation leaves some notation and ambient hypotheses unspecified. The following records its question and makes the usual tracial interpretation of the conclusion explicit, without resolving those omissions.

In the usual tracial framework, let \(M=N_1*_B N_2\) denote an amalgamated free product of finite von Neumann algebras with compatible faithful normal tracial states and trace-preserving conditional expectations onto their common unital von Neumann subalgebra \(B\). Write \(\tau\) for the resulting trace on \(M\), and \(E_B:M\to B\) for the canonical conditional expectation. These ambient hypotheses explain the notation; they are not explicitly stated in the supplied record.

For an index set \(I\) and an ultrafilter \(\mathcal U\) on \(I\), the tracial ultrapower is
\[
M^{\mathcal U}=\ell^\infty(I,M)\big/\{(x_t):\lim_{t\to\mathcal U}\tau(x_t^*x_t)=0\},
\qquad
\tau^{\mathcal U}((x_t)_{\mathcal U})=\lim_{t\to\mathcal U}\tau(x_t).
\]
Identify \(M\) with its diagonal image in \(M^{\mathcal U}\), and let
\[
E_{B^{\mathcal U}}((x_t)_{\mathcal U})=(E_B(x_t))_{\mathcal U}.
\]
The source does not specify the index set or restrictions on \(\mathcal U\).

The symbols \(\mathcal H(A)\) and \(\sim_A\) are not defined in the supplied record and are retained as source notation. If \(\mathcal H(A)\) is intended to mean the Haar unitaries of a tracial von Neumann algebra \((A,\tau_A)\), its definition would be
\[
\mathcal H(A)=\{u\in A:u^*u=uu^*=1,\ \tau_A(u^n)=0\text{ for every }n\in\mathbb Z\setminus\{0\}\}.
\]
This interpretation requires confirmation; no meaning is assigned here to \(\sim_A\).

\textbf{Source question, conditional on resolving the notation.} For every pair \(u_i\in\mathcal H(N_i)\), \(i\in\{1,2\}\), satisfying
\[
E_B(u_i^k)=0\qquad(k\geq1),
\]
and every pair \(v_i\in\mathcal H(M^{\mathcal U})\) satisfying
\[
u_i\sim_{M^{\mathcal U}}v_i,
\qquad E_{B^{\mathcal U}}(v_i^k)=0
\qquad(i\in\{1,2\},\ k\geq1),
\]
must \(v_1\) and \(v_2\) be freely independent?

If the intended conclusion is scalar-valued freeness with respect to \(\tau^{\mathcal U}\), it means the following: with \(A_i=W^*(v_i)\), for every \(m\geq1\), every sequence \(i_1,\ldots,i_m\in\{1,2\}\) with \(i_j\ne i_{j+1}\) for \(1\leq j<m\), and every choice \(a_j\in A_{i_j}\) satisfying \(\tau^{\mathcal U}(a_j)=0\), one has
\[
\tau^{\mathcal U}(a_1\cdots a_m)=0.
\]
The supplied wording does not explicitly distinguish this conclusion from freeness with amalgamation over \(B^{\mathcal U}\). Throughout the displayed moment hypotheses, the source's \(k\in\mathbb N\) is interpreted as \(k\geq1\); including \(k=0\) would contradict unitality of the conditional expectations in a nonzero algebra.
\subsection{Short English statement}
The question asks whether two elements in a tracial ultrapower must be freely independent when each is related, by the source's undefined relation \(\sim\), to an element from a different amalgamated free-product component, and all their positive powers have zero conditional expectation onto the amalgamating algebra. The meanings of \(\mathcal H\), \(\sim\), the ambient hypotheses, and the intended type of freeness require confirmation.
\subsection{Sources}
\begin{itemize}
\item Supplied local UnsolvedMath registry record: permanent ID \(801126\), exact source code \texttt{catalogue entry}, title \emph{Ultraproduct II\_1 Factors --- Problem 27}. The registry attributes the problem to Elayavalli, with proposed year \(2025\), and labels it \texttt{open}; that label is supplied curation, not an independently verified mathematical assessment.
\item Supplied source-list attribution: Elayavalli, \emph{50 Open Problems: Ultraproduct II\_1 Factors} (2025), Problem 27. Supplied URL: \url{https://arxiv.org/abs/2511.20377}. The background records access on \texttt{2026-10-16} and extraction method \texttt{arxiv\_latex}; these are supplied provenance claims, not a claim of independent source access.
\item Supplied upstream revision: \texttt{d371ee917ac9647fbb06cd6fc529deec6f3dda82}; registry creation and update timestamps: \texttt{2026-10-16T14:00:00Z}.
\end{itemize}
\end{document}
